% TOP_PROBLEM: {"problemId": "problem.polynomial-time-classical-integer-factorization", "canonicalTitle": "Polynomial-time classical integer factorization", "releaseRank": 17, "primaryDomain": "theoretical_computer_science", "primaryDomainLabel": "Theoretical computer science", "displayStatus": "open", "releaseStatus": "open"}
% !TeX program = lualatex
% Definition and sources only; this file is not an LLM solution attempt.
\documentclass[11pt]{article}
\usepackage[margin=25mm]{geometry}
\usepackage{fontspec}
\setmainfont{Latin Modern Roman}
\usepackage{amsmath,amssymb,mathtools}
\usepackage{unicode-math}
\setmathfont{Latin Modern Math}
\usepackage{xurl,hyperref}
\hypersetup{colorlinks=true,urlcolor=blue}
\setcounter{secnumdepth}{0}
\setlength{\parindent}{0pt}
\setlength{\parskip}{5pt}
\setlength{\emergencystretch}{3em}
\begin{document}
\section*{\texorpdfstring{Polynomial-time classical integer factorization}{Polynomial-time classical integer factorization}}\label{17}
\subsection{Definitions and mathematical statement}
The binary length of $N\ge2$ is $\ell(N)=\lfloor\log_2N\rfloor+1$. A nontrivial factor is an integer $d$ with $1<d<N$ and $d\mid N$. Ask whether there exist a deterministic classical algorithm $A$ and constants $C,k$ such that
\[
 \forall N\text{ composite},\qquad A(N)=d\text{ a nontrivial factor},\qquad
 T_A(N)\le C\ell(N)^k.
\]
The promise is compositeness; the running-time bound is in the bit length, not in the value of $N$.
\par\smallskip\textit{Formulation and concept references:} \href{https://www.sciencedirect.com/science/article/pii/S0022000015000768}{[S1]}.

\subsection{Short English statement}
Is there a deterministic classical method that factors every composite integer in time polynomial in its number of binary digits?
\subsection{Sources}
\begin{itemize}
\item[S1]Integer factoring and modular square roots.
\url{https://www.sciencedirect.com/science/article/pii/S0022000015000768}.
\textit{Formulation source.}
\item[S2]A number-theoretic conjecture implying faster algorithms for polynomial factorization and integer factorization — Chris Umans and Siki Wang (2025).
\url{https://arxiv.org/abs/2511.10851}.
\textit{Further reference; not a proof endorsement.}
\item[S3]Improved Separation Between Quantum and Classical Computers for Sampling and Functional Tasks — Marshall, Aaronson and Dunjko, CCC 2025.
\url{https://drops.dagstuhl.de/storage/00lipics/lipics-vol339-ccc2025/LIPIcs.CCC.2025.5/LIPIcs.CCC.2025.5.pdf}.
\textit{Further reference; not a proof endorsement.}
\end{itemize}
\end{document}
